\begin{resumo}

    %Este Mapeamento Sistemático da Literatura (MSL) aborda a relevância dos jogos digitais sérios e educacionais na conscientização, prevenção e combate à violência sexual. A pesquisa tem como objetivo identificar e caracterizar os jogos existentes nesse contexto, analisando sua eficácia, impacto e abordagens de desenvolvimento. Utilizando critérios específicos de inclusão e exclusão, o estudo busca responder a questões fundamentais sobre os estilos de narrativa, formatos e plataformas dos jogos, visando compreender como esses elementos influenciam a experiência do usuário e a transmissão de mensagens educacionais. Por meio de uma análise detalhada dos artigos selecionados, este trabalho visa fornecer insights valiosos para pesquisadores, profissionais, educadores e desenvolvedores de jogos, contribuindo para o avanço do conhecimento nessa área e orientando futuras pesquisas e práticas.
    
    A Inteligência Artificial Generativa (IAG) vem se consolidando como um dos principais vetores de transformação na produção de conteúdos audiovisuais não interativos, impulsionada pela maturação técnica de arquiteturas como Modelos de Difusão e Transformers. Apesar do crescimento expressivo das aplicações em fluxos reais de produção, a literatura permanece fragmentada e as propostas taxonômicas existentes tratam de forma isolada as capacidades técnicas, as etapas do \textit{pipeline} ou os riscos éticos, o que impede a análise conjunta dos \textit{trade-offs} implicados em cada escolha tecnológica. Diante desse cenário, esta pesquisa desenvolve a TIGA (\textit{Taxonomy for Integrating Generative AI in Non-Interactive Audiovisual Production}), uma taxonomia multidimensional de integração da IAG no audiovisual construída pelo método iterativo de desenvolvimento taxonômico sobre a base empírica de um Mapeamento Sistemático da Literatura (MSL) prévio, que consolidou sessenta estudos primários recuperados em uma janela de busca de 2022 a abril de 2025, efetivamente publicados entre 2023 e abril de 2025. Dado o volume crescente de publicações e a rápida obsolescência das ferramentas comerciais, o rigor metodológico foi concentrado na síntese qualitativa e na reinterpretação analítica dos dados já extraídos, subordinada à meta-característica \textit{integração da IAG no pipeline de produção audiovisual não interativo}, em vez de novas rodadas quantitativas de busca. A TIGA é facetada: cada instância recebe, simultaneamente, um perfil operacional (grau de autonomia, etapa do \textit{pipeline} e modo de interação), um perfil técnico (família arquitetural, modelo base, tipo de saída, infraestrutura e gargalo dominante) e um perfil de risco e ética (autoria, viés, impacto no trabalho, transparência, reprodutibilidade e perceptibilidade), acompanhado de rubrica de decisão, protocolo de aplicação e mapa de \textit{trade-offs}. A aplicação a três estudos reais do corpus confirmou a correlação entre maior autonomia delegada à máquina e maior custo sociotécnico. Este documento entrega a versão preliminar do artefato e os instrumentos de sua avaliação; a validação plena, por painel de especialistas e por medida de concordância entre avaliadores, constitui a etapa seguinte da pesquisa.
    
    
    \end{resumo}
    